% Standalone editable report. Compile with pdfLaTeX, XeLaTeX, or Tectonic.
% build_report.py refreshes only the two GENERATED blocks from results.csv.
\documentclass[11pt,a4paper]{article}
\usepackage[margin=19mm,headheight=14pt]{geometry}
\usepackage{amsmath,amssymb,booktabs,xcolor,listings,fancyhdr}
\usepackage[hidelinks]{hyperref}
\definecolor{navy}{RGB}{20,43,71}
\definecolor{teal}{RGB}{0,102,112}
\definecolor{codebg}{RGB}{242,247,249}
\setlength{\parindent}{0pt}
\setlength{\parskip}{6pt}
\pagestyle{fancy}
\fancyhf{}
\fancyhead[L]{\small\sffamily IIT TIRUPATI / AI ACCELERATOR DESIGN}
\fancyfoot[L]{\small Assignment 3 / BCHW tensor layout / 6 October 2026}
\fancyfoot[R]{\small\thepage\ / 3}
\renewcommand{\headrulewidth}{0.4pt}
\newcommand{\reporttitle}[1]{{\LARGE\sffamily\bfseries\color{navy}#1}\par\medskip}
\newcommand{\reportsection}[1]{{\large\sffamily\bfseries\color{teal}#1}\par}
\lstset{basicstyle=\ttfamily\small,backgroundcolor=\color{codebg},
  frame=none,xleftmargin=6pt,xrightmargin=6pt,aboveskip=5pt,belowskip=6pt,
  columns=fullflexible,keepspaces=true,showstringspaces=false}
% Fill these directly or pass --student and --roll to build_report.py.
\newcommand{\StudentName}{}
\newcommand{\RollNumber}{}
\newcommand{\EmptyIdentity}{}
% BEGIN GENERATED METADATA
\newcommand{\ResultsSHA}{6d20234babae59337033c2b75f0d99b5e20c15391c590660cfaad1c3c5f8b634}
\newcommand{\ElementCount}{451,492}
\newcommand{\PythonVersion}{3.14.2}
\newcommand{\NumpyVersion}{2.4.2}
% END GENERATED METADATA
\hypersetup{pdftitle={Assignment 3 - BCHW Tensor Layout}}
\begin{document}
\reporttitle{Indexing and reconstruction}
\ifx\StudentName\EmptyIdentity\else\textbf{Name:} \StudentName\quad\fi
\ifx\RollNumber\EmptyIdentity\else\textbf{Roll no.:} \RollNumber\par\fi
\textbf{Objective.} Map $I\in\mathbb{R}^{B\times C\times H\times W}$ to
$F\in\mathbb{R}^{B\times(CHW)}$ and reconstruct $\widehat I$.
The batch index is preserved; width varies fastest, followed by height and channel.
This linear ordering supports tensor buffering and accelerator data transfers.

\reportsection{1. Forward mapping}
\begin{lstlisting}
F = allocate(B, C*H*W, same dtype as I)
for b in 0..B-1, c in 0..C-1:
    for h in 0..H-1, w in 0..W-1:
        k = (c*H + h)*W + w
        F[b,k] = I[b,c,h,w]
\end{lstlisting}
\reportsection{2. Reverse mapping}
\begin{lstlisting}
I_hat = allocate(B, C, H, W, same dtype as F)
for b in 0..B-1, k in 0..C*H*W-1:
    c = k // (H*W)
    r = k % (H*W)
    h = r // W; w = r % W
    I_hat[b,c,h,w] = F[b,k]
\end{lstlisting}
Euclidean division gives a unique decomposition $k=c(HW)+hW+w$, where
$0\leq h<H$ and $0\leq w<W$. The mapping is therefore bijective.
Both conversions take $O(BCHW)$ time and allocate $O(BCHW)$ output storage.
The reverse conversion requires the original $C,H,W$ as metadata.

\reportsection{3. Manual example: $B=1, C=2, H=2, W=3$}
\begin{lstlisting}
I[0,0,:,:] = [[0,1,2], [3,4,5]]
I[0,1,:,:] = [[6,7,8], [9,10,11]]
F[0,:] = [0,1,2,3,4,5,6,7,8,9,10,11]
\end{lstlisting}
\begin{center}
\begin{tabular}{lll}
\toprule
Source $(b,c,h,w)$ & Forward calculation & Destination $(b,k)$ \\
\midrule
$(0,0,0,2)$ & $(0\cdot2+0)\cdot3+2=2$ & $(0,2)$ \\
$(0,0,1,0)$ & $(0\cdot2+1)\cdot3+0=3$ & $(0,3)$ \\
$(0,1,0,0)$ & $(1\cdot2+0)\cdot3+0=6$ & $(0,6)$ \\
$(0,1,1,2)$ & $(1\cdot2+1)\cdot3+2=11$ & $(0,11)$ \\
\bottomrule
\end{tabular}
\end{center}
\textbf{Inverse example.} For $k=10$, $c=\lfloor10/6\rfloor=1$,
$r=10\bmod6=4$, $h=\lfloor4/3\rfloor=1$, and $w=4\bmod3=1$.
Thus $F[0,10]$ returns to $\widehat I[0,1,1,1]=10$.

\newpage
\reporttitle{Experiments and measured errors}
Every run checks all flattened elements, reconstructed elements, and inverse
indices. The signed difference tensor $E=I-\widehat I$ is computed explicitly
and saved in \texttt{errors/00.npy} through \texttt{errors/09.npy} in table order.
The two error measures include all $BCHW$ elements:
\[
E_{\max}=\max_{b,c,h,w}|I[b,c,h,w]-\widehat I[b,c,h,w]|,
\]
\[
\operatorname{MAE}=\frac{1}{BCHW}
\sum_{b,c,h,w}|I[b,c,h,w]-\widehat I[b,c,h,w]|.
\]
\reportsection{4. Results (measured values from results.csv)}
\begin{center}
\renewcommand{\arraystretch}{1.4}
\begin{tabular}{lrrrrrr}
\toprule
Input & $B$ & $C$ & $H$ & $W$ & $E_{\max}$ & MAE \\
\midrule
% BEGIN GENERATED RESULTS
Manual ordering & 1 & 2 & 2 & 3 & 0.0 & 0.0 \\
MNIST test samples 0-1 & 2 & 1 & 28 & 28 & 0.0 & 0.0 \\
Synthetic RGB & 2 & 3 & 32 & 32 & 0.0 & 0.0 \\
Synthetic fmap C=8 & 2 & 8 & 13 & 17 & 0.0 & 0.0 \\
Synthetic fmap C=16 & 2 & 16 & 13 & 17 & 0.0 & 0.0 \\
Synthetic fmap C=32 & 2 & 32 & 13 & 17 & 0.0 & 0.0 \\
Synthetic fmap C=64 & 2 & 64 & 13 & 17 & 0.0 & 0.0 \\
Synthetic fmap C=128 & 2 & 128 & 13 & 17 & 0.0 & 0.0 \\
Synthetic fmap C=256 & 2 & 256 & 13 & 17 & 0.0 & 0.0 \\
Synthetic fmap C=500 & 2 & 500 & 13 & 17 & 0.0 & 0.0 \\
% END GENERATED RESULTS
\bottomrule
\end{tabular}
\end{center}
\textbf{10 / 10 experiments passed; \ElementCount{} elements verified.}

\textbf{MNIST.} The first two test images (indices 0 and 1) are loaded from
the original IDX archive as unnormalized \texttt{uint8} grayscale values.
The archive is included in \texttt{data/}; its SHA-256 checksum and IDX header
are validated before use.

\textbf{Synthetic data.} The three-channel input is synthetic RGB-shaped data,
not CIFAR. Higher-channel inputs are synthetic \texttt{float32} feature maps,
not outputs of a trained network. Each case uses a fresh NumPy PCG64 generator
with seed $20261005+C$ and a standard normal distribution. Dimensions remain
practical: the largest experiment contains 221,000 elements.

All measured maximum absolute errors and MAEs are 0.0, and every saved signed
difference is zero. Neither conversion performs arithmetic on tensor values;
both preserve dtype, so exact reconstruction is expected.

\vfill
{\footnotesize Results CSV SHA-256:\par\scriptsize\texttt{\ResultsSHA}}

\newpage
\reporttitle{Validation and discussion}
\reportsection{5. Verification}
The 16 automated unit tests cover handwritten expected ordering, $B>1$,
non-square height/width, singleton axes, all nine requested channel counts,
four numeric dtypes, noncontiguous inputs, index bijection, input preservation,
invalid inputs, and explicit-loop implementation. Deliberate matrix permutations
and reconstruction corruption must be detected.

An independent test uses unsigned inputs with expected signed differences
$[-255,255,2,-3]$, maximum absolute error 255, and MAE 128.75.
This verifies the error calculator and guards against unsigned subtraction
wraparound. Integer values are subtracted before conversion to floating-point
metrics; a separate test checks differences between large \texttt{int64} values.

The conversion call audit allows only validation, allocation, loops and error
construction. No \texttt{reshape}, \texttt{view}, \texttt{flatten}, \texttt{ravel},
or equivalent library layout operation is used in either algorithm. No optional
library layout-conversion reference is used.

\reportsection{6. Observations and limitations}
Zero error from $C=1$ through $C=500$ supports correct channel boundaries and
preservation of the batch dimension. An independent ordering check is necessary:
two mutually consistent but incorrect permutations could pass a round trip.
Quotient/remainder tests establish that each valid linear address has exactly
one BCHW coordinate.

The implementation uses Python loops to explain address generation, not to
benchmark accelerator throughput. It supports positive dimensions and finite
real integer/floating-point NumPy arrays. Empty axes, NaN, infinity, complex
and object values are rejected. Differences and aggregate metrics use
\texttt{float64} storage; exact integer subtraction precedes storage. Arbitrary
precision error reporting for extreme values is outside these experiments.

\reportsection{7. Reproducibility and sources}
The measured experiments used Python \PythonVersion{} and NumPy \NumpyVersion{}.
Run the unit tests and \path{experiments.py} to reproduce the results.

Next, \path{build_report.py} refreshes measured content in this editable source
and invokes LaTeX to compile its PDF. Run \path{verify_submission.py} to compare
the LaTeX/PDF tables with the CSV, console, manifest and saved differences.
Setup commands, data provenance and the requirement checklist are in
\path{EVALUATOR.md}.

\textbf{Assignment source:} user-provided \texttt{assgn3-tensor-flat.pdf}, pages 1--3.

\textbf{Dataset source:} CVDF's public mirror of the MNIST test images.
\par\url{https://github.com/cvdfoundation/mnist}
\par The archive was downloaded on 5 October 2026 and included unchanged
for offline experiment reproduction.
The RGB-shaped and higher-channel cases are explicitly synthetic.

\textbf{Editing:} this standalone \texttt{report.tex} can be compiled in Overleaf,
with pdfLaTeX/XeLaTeX, or with Tectonic. Edit the name/roll macros near the top.
The generated blocks preserve the measured results; other prose is freely editable.
\end{document}
